\documentclass[10pt,a4paper]{amsart}
\usepackage[margin=19mm]{geometry}
\usepackage{amsmath,amssymb,mathtools}
\usepackage[scr=rsfs,cal=euler]{mathalfa}
\usepackage{xfrac}
\usepackage[unicode,hidelinks,bookmarks=false]{hyperref}
\newcommand{\ie}{i.e.\ }
\newcommand{\cf}{cf.\ }
\newcommand{\resp}{respectively\ }
\newcommand{\Subsection}{\subsection}
\providecommand{\nobreakdash}{\nobreak\hbox{-}}
\setlength{\emergencystretch}{2em}
\multlinegap0pt
\usepackage{bm}
\usepackage{xcolor}
\usepackage{enumerate}
\newtheorem{lemma}{Lemma}
\newcommand{\D}{\mathrm{d}} 
\newcommand{\Ge}{\geqslant}
\DeclareMathOperator{\KL}{KL}
\newcommand{\Le}{\leqslant}
\newcommand{\ep}{\varepsilon}
\title{Jeffreys Flow: Robust Boltzmann Generators for Rare Event Sampling via Parallel Tempering Distillation}
\author{Guang Lin, Christian Moya, Di Qi, Xuda Ye}
\date{Source: arXiv:2604.05303v1 (CC BY 4.0)}
\begin{document}
\maketitle

\noindent\emph{Source and licence.} Jeffreys Flow: Robust Boltzmann Generators for Rare Event Sampling via Parallel Tempering Distillation. Version arXiv:2604.05303v1 (CC BY 4.0), distributed by arXiv under the Creative Commons Attribution 4.0 International licence, http://creativecommons.org/licenses/by/4.0/, as recorded in the arXiv licence metadata for this exact version and shown on the abstract page https://arxiv.org/abs/2604.05303v1. Full licence notice: \texttt{licenses/arxiv-2604.05303-CC-BY-4.0.txt}.\par\smallskip

\noindent\textit{Editorial scope.} Counted unit: the article's lemma concentration (source0.tex lines 239--246). The article's complete original proof of it is delivered with it (source0.tex lines 247--274, one \texttt{proof} environment of 28 lines). No mathematics is written, repaired or re-derived: the statement and the proof are the article's own lines, in the article's own notation, and no retained line is substituted or re-flowed.  No further source-local context is needed: every cross-reference of every retained line resolves inside the two delivered spans.  Nothing was omitted from any retained span, so the package carries no omission record.  No retained line contains a citation, so the wrapper carries no bibliography and no bibliography entry.  No retained line contains a figure environment, an image-inclusion command or drawing code, so no image is delivered and the package declares no asset dependency.  Editorial formatting only, in addition: an AMS \texttt{amsart} preamble that restates the article's own declarations and macros used by the retained lines, namely the environments \texttt{lemma} on separate counters, together with the standard packages those lines need.  The retained environments print in this document as Lemma 1 in order of appearance, whereas the article prints the same items with the numbering of its own full document.  The complete native source of the article as distributed in the arXiv e-print for this exact version is bundled unchanged as \texttt{sources/197948/source0.tex}.
\begin{lemma}
	\label{lemma: concentration}
	Assume that there exists $\epsilon>0$ such that the distributions $P$ and $Q$ satisfy $D_{\KL}(P\|Q) \Le \epsilon$ and $D_{\KL}(Q\|P) \Le \epsilon$. Then the probability of $A_\delta$ is bounded by
	\begin{equation}
		\max\big(P(A_\delta),Q(A_\delta)\big) \Le \frac{2\ep}{\delta(1-e^{-\delta})}.
		\label{eq: mPQb}
	\end{equation}
\end{lemma}

\begin{proof}
By definition, the KL divergence bounds yield
\begin{equation}
\int_{\mathbb R^d} (P(x) - Q(x))\log\frac{P(x)}{Q(x)} \D x \Le 2\epsilon.
\label{eq: PQep}
\end{equation}
Note that the integrand in \eqref{eq: PQep} is non-negative. We partition the instability set $A_\delta$ into
\begin{equation*}
	S_+ = \{P(x) \Ge e^\delta Q(x)\}, \quad S_- = \{Q(x) \Ge e^\delta P(x)\},
\end{equation*}
and define $I_+$ and $I_-$ to be their respective contributions to the integral in \eqref{eq: PQep}. Then \eqref{eq: PQep} implies $I_++I_-\Le 2\epsilon$.

On $S_+$, we have $\log\frac{P(x)}{Q(x)} \Ge \delta$ and
\begin{equation*}
	I_+ \Ge \int_{S_+} P(x)(1-e^{-\delta})\delta \D x = \delta(1-e^{-\delta})P(S_+).
\end{equation*}
Since $Q(x) \Le P(x)$ on $S_+$, we find
\begin{equation}
	Q(S_+) \Le P(S_+) \Le \frac{I_+}{\delta(1-e^{-\delta})}.
	\label{eq: QPd}
\end{equation}
By symmetry, on $S_-$ we obtain
\begin{equation}
	P(S_-) \Le Q(S_-) \Le \frac{I_-}{\delta(1-e^{-\delta})}.
	\label{eq: PQd}
\end{equation}
Since $I_++I_-\Le 2\ep$, adding \eqref{eq: QPd} and \eqref{eq: PQd} yields \eqref{eq: mPQb}.
\end{proof}

\end{document}
